\documentclass[10pt,letterpaper]{article}

\usepackage[letterpaper,top=0.45in,bottom=0.45in,left=0.48in,right=0.48in]{geometry}
\usepackage[T1]{fontenc}
\usepackage{newpxtext,newpxmath}
\usepackage{microtype}
\usepackage{enumitem}
\usepackage{titlesec}
\usepackage{xcolor}
\usepackage[hidelinks]{hyperref}
\usepackage{tabularx}

% NOTE FOR KESHAV: every \est{...} value is an estimate, not a measured
% number. Cross-check each one against metrics-to-verify.md, replace it with
% the real value, then redefine \est as \newcommand{\est}[1]{#1} (or delete
% the wrappers) to render everything in black again.
\newcommand{\est}[1]{\textcolor{red!70!black}{#1}}

\hypersetup{pdftitle={Keshav Krishna - Resume}}

\pagestyle{empty}
\setlength{\parindent}{0pt}
\setlength{\tabcolsep}{0pt}
\setlength{\emergencystretch}{2em}
\urlstyle{same}

\titleformat{\section}
  {\large\scshape}
  {}{0pt}{}
  [\vspace{0.12em}\titlerule]
\titlespacing*{\section}{0pt}{1.2em}{0.65em}

\setlist[itemize]{
  label=--,
  leftmargin=1.8em,
  itemsep=0.3em,
  topsep=0.3em,
  parsep=0pt,
  partopsep=0pt
}

\newcommand{\entry}[4]{%
  \begin{tabularx}{\textwidth}{@{}Xr@{}}
    \textbf{#1} & #2 \\
    \textit{#3} & \textit{#4}
  \end{tabularx}\vspace{-0.2em}
}

\newcommand{\project}[2]{%
  \begin{tabularx}{\textwidth}{@{}Xr@{}}
    \textbf{#1} & \textit{#2}
  \end{tabularx}\vspace{-0.25em}
}

\begin{document}
\fontsize{10.4}{13.2}\selectfont

\begin{center}
  {\fontsize{24}{27}\selectfont\scshape Keshav Krishna}\\[0.15em]
  New York City, US $\mid$ +1 (479) 685 5463 $\mid$ \href{mailto:kk6081@nyu.edu}{kk6081@nyu.edu}\\[-0.05em]
  \href{https://github.com/k-keshav-k}{GitHub} $\mid$
  \href{https://k-keshav-k.github.io/}{Website} $\mid$
  \href{https://www.linkedin.com/in/keshav-krishna-5054681a9/}{LinkedIn} $\mid$
  \href{https://scholar.google.com/citations?user=zN7P_v4AAAAJ}{Google Scholar}
\end{center}

\vspace{0.35em}
Product engineer with three years shipping ML, NLP, and foundation-model systems to production, including consumer systems serving 100M+ users. Currently building risk analytics tooling at Moore Capital Management after extending a full-time summer internship into a part-time fall role. I own features from an ambiguous conversation through production, and I measure model quality with explicit evaluation and benchmarking frameworks.

\section{Education}

\entry{New York University, Courant Institute --- MS, Computer Science, Artificial Intelligence}{Aug 2025 -- May 2027}{GPA 3.95/4.0; Machine Learning, Building LLM Reasoners, Honors Algorithms}{New York City}

\vspace{0.45em}
\entry{Indian Institute of Technology, Ropar --- BTech, Computer Science and Engineering}{Aug 2019 -- Jun 2023}{GPA 8.74/10; Merit Scholarship for top academic performance}{Ropar, India}

\vspace{0.65em}
\section{Experience}

\entry{Moore Capital Management --- Risk Technology Intern}{May 2026 -- Present}{Risk analytics tooling and dashboards for internal users; full-time summer, part-time fall}{New York City}
\begin{itemize}
  \item Owned risk tooling end to end: Python data workflows behind HTML/CSS dashboards, gone from verbal stakeholder request to production in days and now used daily by \est{20+} risk, portfolio, and operations staff.
  \item Optimized VaR, CVaR, beta, and correlation calculations with vectorized Python and Numba JIT, cutting runtime about \est{70\%} on multi-million-row datasets so daily risk numbers are ready before market open.
  \item Benchmarked distributed data-processing tools on the firm's compute cluster with profiling and system modeling, locating scalability bottlenecks that guided infrastructure decisions and lifting batch throughput about \est{2x}.
\end{itemize}

\vspace{0.65em}
\entry{JioSaavn --- Software Engineer, Data Science}{Jul 2023 -- Jul 2025}{Production ML and NLP in a consumer app serving 100M+ users}{Mumbai, India}
\begin{itemize}
  \item Built recommendation and personalization systems on embeddings and vector retrieval (Spark, Hive, Redis, Qdrant), driving \textbf{1M impressions, 800K clicks, and 500K streams per day} in production.
  \item Shipped a LLaMA3-8B-Instruct foundation-model system that combines intent understanding, retrieval-augmented generation, and automated song selection; owned deployment and inference optimization on Docker and Kubernetes, serving \est{50K+} generated playlists per day at sub-\est{2s} p95 latency.
  \item Designed the benchmark dataset, evaluation metrics, and acceptance criteria for multilingual playlist moderation, then ran a bake-off across SetFit, HingRoBERTa, mURIL, and LLaMA3-8B-Instruct; the selected model hit \textbf{85\% accuracy} and absorbed about \est{60\%} of the manual review load.
  \item Owned model, pipeline, and deployment for \est{3+} shipped systems over two years, working cross-functionally with product managers, domain experts, and platform engineering teams.
\end{itemize}

\vspace{0.65em}
\entry{GE Healthcare --- Software Intern, Medical Computer Vision}{May 2022 -- Jul 2022}{Detection and de-identification pipelines for clinical imaging}{Bangalore, India}
\begin{itemize}
  \item Built YOLO-based medical-image de-identification and RGB/IR people-detection pipelines over 20K+ images, reaching \est{95\%} detection accuracy and cutting manual image review about \est{80\%}.
\end{itemize}

\section{Selected Technical Projects}

\project{FlashAttention-2 --- from-scratch Triton implementation~\href{https://github.com/k-keshav-k/Flash-Attention}{[GitHub]}}{2026}
\begin{itemize}
  \item Implemented forward and backward kernels with online softmax, tiling, causal masking, and activation recomputation; cut attention memory from $O(N^2)$ to $O(N)$, taking 65K-token attention from \est{10+ GB} to \est{under 1 GB} and making long-context training feasible.
  \item Profiled Transformer models up to 2.7B parameters with NVIDIA Nsight Systems and PyTorch memory snapshots, benchmarking against vanilla PyTorch and \texttt{torch.compile}; measured about \est{2x} speedup over naive attention.
\end{itemize}

\vspace{0.6em}
\project{Transformer language model from scratch~\href{https://github.com/k-keshav-k/Large-Language-Models-from-Scratch}{[GitHub]}}{2026}
\begin{itemize}
  \item Built a 17M-parameter model in PyTorch from low-level components (BPE tokenization, RoPE, RMSNorm, SwiGLU, causal attention, AdamW, cosine LR scheduling), reaching validation loss \est{3.1}, with multiprocessing tokenization and memory-mapped loading cutting data-pipeline time about \est{3x}.
  \item Ran systematic ablations across pre- vs. post-norm, RoPE vs. NoPE, and SwiGLU vs. SiLU; the pre-norm + RoPE + SwiGLU recipe won, cutting validation loss about \est{8\%} against the baseline configuration.
\end{itemize}

\vspace{0.6em}
\project{Pico-LLM --- training, evaluation, and inference platform~\href{https://github.com/k-keshav-k/pico-llm}{[GitHub]}}{2025}
\begin{itemize}
  \item Built a unified platform for SFT and DPO/GRPO training, Model Ops (checkpoint management, batch evaluation), and inference across Transformer, LSTM, and K-gram architectures.
  \item Added KV-cached generation and multi-GPU batch generation, lifting inference throughput about \est{3x} and improving experiment reliability across architectures.
\end{itemize}

\vspace{0.6em}
\project{Alignment and Reasoning RL for Mathematical LLMs~\href{https://github.com/k-keshav-k/Alignment-and-Reasoning-RL}{[GitHub]}}{Spring 2026}
\begin{itemize}
  \item Built an end-to-end reasoning and post-training pipeline: zero-shot evaluation, SFT, GRPO with verified rewards, masked losses, entropy tracking, and policy-gradient updates.
  \item Fine-tuned Qwen2.5-Math-1.5B on chain-of-thought data through SFT and GRPO, lifting MATH accuracy from \textbf{56.4\%} to \textbf{60.4\%} (+4.0 points).
\end{itemize}

\vspace{0.65em}
\section{Research and Publications}

\project{Phase-Adaptive Generation: Efficient Decoding for Diffusion Language Models~\href{https://github.com/k-keshav-k/Phase-Assisted-Generation}{[GitHub]}}{Feb 2026 -- May 2026}
\begin{itemize}
  \item NYU project under Prof. Greg Durrett. Analyzed generation dynamics of Dream/LLaDA-style diffusion LMs via per-token refinement traces, revealing interpretable phase structure: template spans stabilize immediately while arithmetic and answer-finalization spans demand substantially more refinement.
  \item Trained Transformer predictors over 138K block-level control tuples to forecast per-block refinement budgets; integrated into the LLaDA decoding loop, matching AdaBlock's \textbf{89.5\% GSM8K accuracy} with approximately \textbf{21\% fewer forward evaluations}. \href{https://github.com/k-keshav-k/Phase-Assisted-Generation/blob/main/writeup/final_report.pdf}{[Report]}
\end{itemize}

\vspace{0.75em}
\project{ML-Driven Program Phase Prediction for Dynamic Hardware Resource Optimization~\href{https://github.com/k-keshav-k/ML-based-hardware-optimization-using-Phases}{[GitHub]}}{Jan 2026 -- May 2026}
\begin{itemize}
  \item Independent study with Prof. Mohamed Zahran, NYU Courant. Clustered temporal traces of multicore workloads into interpretable pressure states and distilled them into compact predictors with sub-KB model state, reaching \est{90\%+} phase-prediction accuracy while keeping per-timestep overhead negligible.
\end{itemize}

\vspace{0.75em}
\project{Publications}{2024 -- 2025}
\begin{itemize}
  \item \href{https://www.frontiersin.org/journals/artificial-intelligence/articles/10.3389/frai.2025.1441250/full}{\textit{Advancements in Cache Management: A Review of Machine Learning Innovations}} --- Frontiers in AI, 2025. Single-author, 10+ citations; reviewed 40+ papers on ML for cache management.
  \item \href{https://arxiv.org/abs/2410.15344}{\textit{LLC Intra-set Write Balancing for Non-Volatile Memory Caches}} --- arXiv, 2024. Bachelor's thesis research.
\end{itemize}

\section{Technical Skills}

\textbf{Product / full-stack:} React, JavaScript, HTML/CSS, FastAPI, Postgres, Redis, Docker, Kubernetes, Git, Linux, developer enablement and technical documentation\\
\textbf{LLMs and GenAI:} foundation model integration and deployment, production-ready GenAI capabilities, scalable service integration, Model Ops (model hubs, versioning, checkpointing), agentic systems and multi-agent orchestration, tool-use orchestration, reasoning pipelines, agent evaluation frameworks, evaluation and benchmarking pipelines, continuous evaluation and automated model testing, retrieval-augmented generation, vector search (Qdrant), context management, compound AI systems, SFT/GRPO/RLHF, PyTorch, Hugging Face, tokenizers, vLLM, inference optimization, Anthropic and OpenAI APIs, MCP\\
\textbf{Data / infrastructure:} Python, C/C++, Spark, Hive, Airflow, Triton kernel development, NVIDIA Nsight Systems, Slurm/HPC, cloud AI infrastructure (GCP, Azure), reproducible pipelines

\end{document}
